% problem-000169 1 铝铁合金反应方程式
\section*{1. 铝铁合金反应方程式}
\textit{下列为分问。}

\subsection*{1-1}
为获得性能良好的纳⽶材料，利⽤团簇化合物 $\mathrm { K _ { 1 2 } S i _ { 1 7 } }$ 和 $\mathrm { N H _ { 4 } B r }$ 反应制备纳⽶硅。

\subsection*{1-2}
硒代硫酸根在酸性条件下被⼀定浓度的 $\mathrm { H _ { 2 } O _ { 2 } }$ 氧化，主产物与硫代硫酸根被 $\mathrm { I } _ { 2 }$ 氧化的产物相似。

\subsection*{1-3}
单质碲(Te)和 $\mathrm { A s F _ { 5 } }$ 在溶剂 $\mathrm { S O } _ { 2 }$ 中按计量⽐6:3反应，主产物为灰⾊抗磁性固体，其阴离⼦为⼋⾯体结构。

\subsection*{1-4}
利⽤ $\mathrm { X e F _ { 2 } }$ 和溴酸根溶液反应制备⾼溴酸根。

\subsection*{1-5}
钼酸钠 $\left( \mathrm { N a _ { 2 } M o O _ { 4 } } \right)$ 和硫代⼄酰胺 $\mathrm { ( C H _ { 3 } C S N H _ { 2 } ) }$ )混合溶液在⽔热条件下强酸性介质中发⽣反应，制备⼆维材料 $\mathrm { M o S } _ { 2 }$ 。（提⽰：分离 $\mathrm { M o S _ { 2 } }$ 后的酸性溶液中加⼊ $\mathrm { B a C l _ { 2 } }$ ，出现⽩⾊沉淀。）
